% \documentclass[a4paper,12pt]{article}
\documentclass[a4paper,11pt]{article}
% -------------------------------------------------
% Set up the page margins.
% -------------------------------------------------
% \usepackage[text={6.5in,9.5in},centering]{geometry}
% \setlength{\topmargin}{-0.25in}
\usepackage[margin=1in]{geometry}

% -------------------------------------------------
% Load Packages here
% -------------------------------------------------
%\usepackage{hyperref}
\usepackage{graphicx,url,subfig}
\RequirePackage[OT1]{fontenc}
\RequirePackage{amsthm,amsmath,amssymb,amscd}
\RequirePackage[numbers]{natbib}
\RequirePackage[colorlinks,citecolor=blue,urlcolor=blue]{hyperref}
\usepackage{graphics}
\usepackage{enumerate}
\usepackage{datetime}
\usepackage{etex}
% \usepackage[notcite,notref,color]{showkeys}
\usepackage[normalem]{ulem} % either use this (simple) or
\usepackage{soul} % use this (many fancier options)
\makeatother
\numberwithin{equation}{section}
\allowdisplaybreaks
\usepackage{mathrsfs}


% -------------------------------------------------
% check references
% -------------------------------------------------
% \usepackage{refcheck}
% \usepackage[notref,notcite]{showkeys}
% \usepackage[notcite]{showkeys}
% \usepackage{showkeys}

% -------------------------------------------------
% Environments here
% -------------------------------------------------
\theoremstyle{plain}
\newtheorem{theorem}{Theorem}[section]
\newtheorem{corollary}[theorem]{Corollary}
\newtheorem{lemma}[theorem]{Lemma}
\newtheorem{proposition}[theorem]{Proposition}
\newtheorem{exercise}{Exercise}
\theoremstyle{definition}
\newtheorem{definition}[theorem]{Definition}
\newtheorem{hypothesis}[theorem]{Hypothesis}
\newtheorem{assumption}[theorem]{Assumption}

\newtheorem{remark}[theorem]{Remark}
\newtheorem{conjecture}[theorem]{Conjecture}
\newtheorem{example}[theorem]{Example}

% -------------------------------------------------
%  Define short hand symbols.
% -------------------------------------------------
\newcommand{\E}{\mathbb{E}}
\newcommand{\W}{\dot{W}}
\newcommand{\B}{\textbf{B}}
\newcommand{\D}{\mathbb{D}}
\newcommand{\ud}{\ensuremath{\mathrm{d}}}
\newcommand{\Ceil}[1]{\left\lceil #1 \right\rceil}
\newcommand{\Floor}[1]{\left\lfloor #1 \right\rfloor}
\newcommand{\sgn}{\text{sgn}}
\newcommand{\Lad}{\text{L}_{\text{ad}}^2}
\newcommand{\SI}[1]{\mathcal{I}\left[#1 \right]}
\newcommand{\SIB}[2]{\mathcal{I}_{#2}\left[#1 \right]}
\newcommand{\Indt}[1]{1_{\left\{#1 \right\}}}
\newcommand{\LadInPrd}[1]{\left\langle #1 \right\rangle_{\text{L}_\text{ad}^2}}
\newcommand{\LadNorm}[1]{\left|\left|  #1 \right|\right|_{\text{L}_\text{ad}^2}}
\newcommand{\Norm}[1]{\left\|  #1   \right\|}
\newcommand{\Ito}{It\^{o} }
\newcommand{\Itos}{It\^{o}'s }
\newcommand{\spt}[1]{\text{supp}\left(#1\right)}
\newcommand{\InPrd}[1]{\left\langle #1 \right\rangle}
\newcommand{\mr}{\textbf{r}}
\newcommand{\Ei}{\text{Ei}}
\newcommand{\arctanh}{\operatorname{arctanh}}
\newcommand{\ind}[1]{\mathbb{I}_{\left\{ {#1} \right\} }}

\DeclareMathOperator{\esssup}{\ensuremath{ess\,sup}}

\newcommand{\steps}[1]{\vskip 0.3cm \textbf{#1}}
\newcommand{\calB}{\mathcal{B}}
\newcommand{\calD}{\mathcal{D}}
\newcommand{\calE}{\mathcal{E}}
\newcommand{\calF}{\mathcal{F}}
\newcommand{\calG}{\mathcal{G}}
\newcommand{\calK}{\mathcal{K}}
\newcommand{\calH}{\mathcal{H}}
\newcommand{\calI}{\mathcal{I}}
\newcommand{\calL}{\mathcal{L}}
\newcommand{\calM}{\mathcal{M}}
\newcommand{\calN}{\mathcal{N}}
\newcommand{\calO}{\mathcal{O}}
\newcommand{\calT}{\mathcal{T}}
\newcommand{\calP}{\mathcal{P}}
\newcommand{\calR}{\mathcal{R}}
\newcommand{\calS}{\mathcal{S}}
\newcommand{\bbC}{\mathbb{C}}
\newcommand{\bbN}{\mathbb{N}}
\newcommand{\bbP}{\mathbb{P}}
\newcommand{\bbZ}{\mathbb{Z}}
\newcommand{\myVec}[1]{\overrightarrow{#1}}
\newcommand{\sincos}{\begin{array}{c} \cos \\ \sin \end{array}\!\!}
\newcommand{\CvBc}[1]{\left\{\:#1\:\right\}}
\newcommand*{\one}{{{\rm 1\mkern-1.5mu}\!{\rm I}}}
\newcommand{\RR}{\mathbb{R}}
\newcommand{\R}{\mathbb{R}}
\newcommand{\myEnd}{\hfill$\square$}
\newcommand{\ds}{\displaystyle}
\newcommand{\Shi}{\text{Shi}}
\newcommand{\Chi}{\text{Chi}}
\newcommand{\Daw}{\ensuremath{\mathrm{daw}}}
\newcommand{\Erf}{\ensuremath{\mathrm{erf}}}
\newcommand{\Erfi}{\ensuremath{\mathrm{erfi}}}
\newcommand{\Erfc}{\ensuremath{\mathrm{erfc}}}
\newcommand{\He}{\ensuremath{\mathrm{He}}}

\def\crtL{\theta}

\newcommand{\conRef}[1]{(#1)}

\DeclareMathOperator{\Lip}{\mathit{L}}
\DeclareMathOperator{\LIP}{Lip}
\DeclareMathOperator{\lip}{\mathit{l}}

\DeclareMathOperator{\Vip}{\overline{\varsigma}}
\DeclareMathOperator{\vip}{\underline{\varsigma}}
\DeclareMathOperator{\vv}{\varsigma}

% % Set up mdframe for questions.
\newcounter{NQ}
\newcounter{LQ}
\newcounter{ZQ}
\usepackage{xcolor}
\usepackage[framemethod=tikz]{mdframed}
\mdfsetup{%
	middlelinecolor=lightgray,
	middlelinewidth=2pt,
	backgroundcolor=red!20,
	roundcorner=10pt}
\newenvironment{NickQuestion}[1][] %
{\refstepcounter{NQ}
	\begin{mdframed}[backgroundcolor=lightgray]\textbf{Nick's Question \theNQ \: #1~~}~~~\noindent}%
	{\end{mdframed}}
\newenvironment{LeQuestion}[1][] %
{\refstepcounter{LQ}
	\begin{mdframed}[backgroundcolor=red!10]\textbf{Le's Question \theLQ \: #1~~}~~~\noindent}%
	{\end{mdframed}}
\newenvironment{ZhijianQuestion}[1][] %
{\refstepcounter{ZQ}
	\begin{mdframed}[backgroundcolor=blue!10]\textbf{Zhijian's Question \theZQ \: #1~~}~~~\noindent}%
	{\end{mdframed}}
\numberwithin{NQ}{section}
\numberwithin{LQ}{section}
\numberwithin{ZQ}{section}

\usepackage[notref,notcite]{showkeys}

\title{Nick Eisenberg's Meeting Notes}
\author{Le Chen, Nick Eisenberg, Zhijian Wu}
\date{October 2019}

\usepackage{natbib}
\usepackage{graphicx}


\begin{document}

\section{Chapter 1}

\begin{NickQuestion}
	Proof of lemma 1.8.4

	\textcolor{magenta}{The proof is very nice.}
\end{NickQuestion}

\section{Chapter 2}
\begin{NickQuestion}
	In the proof of Theorem 2.4.3 on the bottom of page 34 there are some claims made that I am unable to verify why they are true. First,
	we are considering
	\[
	X_t = X_0 + \int_0^t u(s)dB_s
	\]
	and here $u(s)$ only satisfies
	\[
	\mathbb P \left(\int_0^t u(s)^s ds < \infty \right) = 1
	\]
	so, as in the Kuo text, this stochastic integral exists and a limit in probability of a sequence of stochastic integrals $\int_0^t u_n(s)dB_s$ where the $u_n$ are simple and are in $L^2_{ad}$ and hence $\int_0^t u_n(s)dB_s$ is a martingale and we know the first and second moment. But for the limit $ \int_0^t u(s)dB_s$ we know nothing about the moments. But in this text, it is claimed that
	\[
	\E \left[\int_{t_i}^{t_{i+1}}u^2(s)ds - \left(\int_{t_i}^{t_{i+1}}u(s)dB_s\right)^2 \bigg| \mathcal F_i\right] =0
	\]
	and I have never seen this claimed in any other text.

	\textcolor{magenta}{Note that $u$ is bounded (we are assuming localization has been done). Hence, all moments exists.}
\end{NickQuestion}


\begin{NickQuestion}
	In the proof of theorem 2.5.1, why is it true that
	\[
		E \left(\int_0^t |u_n'(B_s)-1_{(x,\infty)}(B_s)|^2 ds\right) \le \int_0^t \mathbb P \left( |B_s-x|\le \frac{2}{n}\right)ds
	\]
	\textcolor{magenta}{Please explore the mollifier $f$ and $f_n$. Plot these $u_n'$ and see how could it differ from $1_{(x,\infty)}$.}
\end{NickQuestion}

\begin{NickQuestion}
	What is the motivation on page 41 behind the following definition of the backward ito integral
	\[
	\int_0^T u(s)\hat{d}B_s = \lim_{n \to \infty} \sum_{i=0}^{n-2}\left( \frac{n}{T}\int_{(j+1)T/n}^{(j+2)T/n}u(s)ds \right)(B_{(j+1)T/n}-B_{jT/n})
	\]
	The integral in the parenthesis is the average of $u_s$ on $\big((j+1)T/n,(j+2)T/n\big)$.

	\bigskip
	\textcolor{magenta}{Let's skip this part. I know nothing on the backward SDE and not particular interested either.
	It has applications in finance. If you don't do finance, I don't see the point here.}
\end{NickQuestion}

\begin{NickQuestion}[(Problem 2.7)]
	This problem follows easily from Levy's Characterization of Brownian Motion. But I can't figure out how to do this problem with the hint given.

	\bigskip
	\textcolor{magenta}{This is a nice problem. The hints is to ask you to find the distribution of $X_t-X_s$
	through the characteristic function. Once apply Ito's formula, take expectation and remove the martingale part, you should
	arrive at an integral equation. Take derivatives, you arrive at a differential equation. The solution for $a_s\left(t\right):=\E\left(e^{i\lambda \left(X_t-X_s\right)}\right)$
	is
	\begin{align*}
		a_s(t) = e^{-\frac{\lambda^2}{2}(t-s)}.
	\end{align*}
	}
\end{NickQuestion}

\begin{NickQuestion}[(Problem 2.15)]
	First we consider the $2-d$ Ito process
	\[
		\begin{pmatrix}
			M_s
			\\ N_s
		\end{pmatrix}
		=
		\begin{pmatrix}
		E(X^p)
		\\ E(Y^{q'})
		\end{pmatrix}
		+
	\int_0^t {	\begin{pmatrix}
		\varphi_s & 0
		\\0 & \psi_s
		\end{pmatrix}
		\begin{pmatrix}
		\hat{d B_s}
		\\ d B_s
		\end{pmatrix} }
	\]
	and then following the hint we apply Ito's formula to $f(x,y) = x^{1/p}y^{1/q'}$. We get the following derivatives:
	\begin{itemize}
		\item $f_x = \frac{1}{p}x^{\frac{1}{p}-1}y^{\frac{1}{q'}}$
		\item $f_{xx} = \frac{1}{p} \left(\frac{1}{p}-1\right)x^{\frac{1}{p}-2}y^{\frac{1}{q'}}$
		\item $f_y = \frac{1}{q'}x^{\frac{1}{p}}y^{\frac{1}{q'}-1}$
		\item $f_{yy} =\frac{1}{q'} \left(\frac{1}{q'}-1\right)x^{\frac{1}{p}}y^{\frac{1}{q'}-2}$
	\end{itemize}
	Now putting this into Ito's Formula gives us
	\begin{align*}
		M_t^{1/p}N_t^{1/q'}&=\Norm{X}_p\Norm{Y}_{q'} + \frac{1}{p}\int_0^t M_s^{\frac{1}{p}-1}N_s^{\frac{1}{q'}} \varphi_s \: d\hat{B_s}+ \frac{1}{q'}\int_0^t M_s^{\frac{1}{p}}N_s^{\frac{1}{q'}-1}\psi_s \: dB_s
		\\& + \frac{1}{2}\frac{1}{p} \left(\frac{1}{p}-1\right)\int_0^t M_s^{\frac{1}{p}-2}N_s^{\frac{1}{q'}}\varphi_s^2 \: ds + \frac{1}{2}\frac{1}{q'} \left(\frac{1}{q'}-1\right)\int_0^t M_s^{\frac{1}{p}}N_s^{\frac{1}{q'}-2} \psi_s^2 \: ds
	\end{align*}

	I am not really sure where to continue from here. I see that if we let $t \to \infty$ then $M_t^{1/p}N_t^{1/q'} \to XY$ so if we took the expected value of both sides then we would partially have what we are looking for since $E(XY)$ is on the left and $\Norm{X}_p\Norm{Y}_{q'}$ is on the right. But I am not sure about the other terms. Im guessing the Ito integrals in the right have $0$ expectation but I am not sure if that is in fact true since I dont see how it is easily seen that the integrands are in $L^2(\mathcal P)$.

	What I did next was see what $\hat{B_s}$ was equal too and for this we define another Ito Process
		\[
	\begin{pmatrix}
	B_s
	\\ \tilde{B_s}
	\end{pmatrix}
	=
	\int_0^t { \begin{pmatrix}
	1 & 0
	\\0 & 1
	\end{pmatrix}
	\begin{pmatrix}
	d B_s
	\\ \tilde{d B_s}
	\end{pmatrix} }
	\]
	and applied Itos formula to $f(t,x,y) = e^{-t}x+ \sqrt{1-e^{-2t}}y$ which gives us
	\begin{align*}
	\hat{B}_s = e^{-t}B_s+ \sqrt{1-e^{-2t}}\tilde{B}_s = \int_0^t \frac{e^{-2s}}{\sqrt{1-e^{-2s}}}\tilde{B}_s - e^{-s}B_s \: ds + \int_0^t e^{-s} \: dB_s + \int_0^t \sqrt{1-e^{-2s}} \: d\tilde{B}_s
	\end{align*}
	However I am not sure if this step is necessary.

\end{NickQuestion}

\bigskip
\textcolor{magenta}{The followings are some of my thoughts:}
First of all, we need the absolute values:
\[
	\begin{pmatrix}
		M_s
		\\ N_s
	\end{pmatrix}
	=
	\begin{pmatrix}
	\E(|X|^p)
	\\ \E(|Y|^{q'})
	\end{pmatrix}
	+
	\int_0^t {	\begin{pmatrix}
	\varphi_s & 0
	\\0 & \psi_s
	\end{pmatrix}
	\begin{pmatrix}
	\hat{d B_s}
	\\ d B_s
	\end{pmatrix} }.
\]
Second, you forget
\begin{align*}
	f_{xy} = \frac{1}{p} \frac{1}{q'} x^{-1/(p')} y^{-1/q}.
\end{align*}
Hence, by Ito's formula, we will have
\begin{align*}
	M_t^{1/p}N_t^{1/q'}
	= &\Norm{X}_p\Norm{Y}_{q'}
	+ \frac{1}{p}\int_0^t M_s^{\frac{1}{p}-1}N_s^{\frac{1}{q'}} \varphi_s \: d\hat{B_s}
	+ \frac{1}{q'}\int_0^t M_s^{\frac{1}{p}}N_s^{\frac{1}{q'}-1}\psi_s \: dB_s \\
	& + \frac{1}{2}\frac{1}{p} \left(\frac{1}{p}-1\right)\int_0^t M_s^{\frac{1}{p}-2}N_s^{\frac{1}{q'}}\varphi_s^2 \: ds
	+ \frac{1}{2}\frac{1}{q'} \left(\frac{1}{q'}-1\right)\int_0^t M_s^{\frac{1}{p}}N_s^{\frac{1}{q'}-2} \psi_s^2 \: ds\\
	&+\frac{1}{pq'} \int_0^t M_s^{-\frac{1}{p'}} N_s^{-\frac{1}{q}} \psi_s \phi_s d\InPrd{B,\hat{B}}_s\\
	=: & \Norm{X}_p\Norm{Y}_{q'} + \sum_{i=1}^5 I_i.
\end{align*}
Upon taking expectation on both sides, we see that
\begin{align*}
	LHS = \E\left(|X_t Y_t|\right)\ge \E\left(X_tY_t\right).
\end{align*}
We only need to prove that
\begin{align*}
	\E\left(\sum_{i=1}^5 I_i\right)\le 0,
\end{align*}
which is clear for $i=1,2$, in which case we have the expectation equal to zero.
Now for the rest three parts, we see that
\begin{align*}
	I_3+I_4+I_5
	=& -\frac{1}{2pq'} M_s^{\frac{1}{p}-2}N_s^{\frac{1}{q'}-2} \left(
		N_s^2 \varphi_s^2 \frac{q'}{p'} -2M_s N_s \phi_s \psi_s e^{-s} + M_s \psi_s^2 \frac{p}{q}
	\right),
\end{align*}
where we have used the fact that $d\InPrd{B,\hat{B}}_s=e^{-s}ds$.
Hence, the condition become the condition that the quadratic form has no solution:
\begin{align*}
	\Delta = 4 e^{-2s} -4\frac{q'p}{p'q} = 4\left(e^{-2s} - \frac{p-1}{q-1}\right)\le 0
\end{align*}
which is equivalent to
\begin{align*}
	q\le e^{2s}(p-1) +1.
\end{align*}
Hence, the largest $q$ that makes the inequality work is $q=e^{2s}(p-1) +1$.


\section{Chapter 3}
\subsection{Reading Material}
\begin{NickQuestion}[(Smooth Cylindrical Random Variable)]
	Following chapter 3 of the Kuo text, denote C the Wiener space of continuous functions from [0,1] to $\R$, $\mathcal R$ the collection cylindrical sets
	\[
	A = \{\omega \in C : (\omega(t_1),...,\omega(t_n) \in U \big| t_i \in [0,1], n \in \mathbb N, U \in \mathscr B (\R^n)  \}
	\]
	and $\mu$ the associated Wiener measure. Consider the coordinate processes $B(t,\omega)=\omega(t)$ which are Brownian motions under $(C,\mathcal B(\mathcal R), \mu)$. There is a theorem that can be found \href{https://math.stackexchange.com/questions/750049/density-of-cylindrical-random-variables-in-classical-wiener-space}{here}.

	\textbf{Theorem} : The random variables
	\[
	\{f(B(t_1,\cdot),...,B(t_n,\cdot)) \big| n\in \mathbb N,\; f \in \mathscr S (\R^n)  \}
	\]
	are dense in $L^2(C,\mu)$ where $\mathscr S (\R^n)$ are the Schwartz functions.

	How does this relate to $\mathcal S$ defined on page 51 of the intro to Malliavin textbook? How are the cylindrical random variables defined in this case? Why is $\mathcal S$ dense in $L^2(\Omega)$? Even in Diffusions Markov Processes and Martingales - Roger and Williams, they define the cylindrical sigma field in the same way Kuo does- through the use of coordinate processes.
\end{NickQuestion}
\begin{proof}[Nick Solution]
	\begin{theorem} \label{density}
		Consider the probability space $(\Omega, \mathcal F, P)$ where $\mathcal F$ is the standard Brownian filtration. Then $\{f(B(h_1),...,f(B(h_n)), n\in \mathbb N,\; f\in C^\infty_p(\R^n), h_i \in H:=L^2(\R^+)   \}$ is dense in $L^2(\Omega)$.
	\end{theorem}
\textbf{Proof:} This proof follows the same scheme as the solution found \href{https://math.stackexchange.com/questions/750049/density-of-cylindrical-random-variables-in-classical-wiener-space}{here} which utilizes Dynkin's multiplicative system theorem. First, we can easily show that the sigma field generated by $\{B(h)\}_{h \in H}$ generates $\mathcal F$. This is true since $1_{[0,t]} \in H$ and $B(1_{[0,t]})=B_t$. Second, since $\mathscr S (\R) \subset C^\infty_p(\R)$ and since for every $h \in H$ there exists functions $f_n \in \mathscr S(\R)$ such that $f_n \to h$ in $L^2(\R)$ then by the Wiener Isometry we have $E|\int_0^\infty h(s)-f_n(s) dB_s|^2 = \int_0^\infty|h(s)-f_n(s)|^2 \ud s \to 0$. This implies that the sigma field generated by $\{f(B(h))\}$ where $h \in H$ and $f \in C^\infty_p(\R)$ is the same as $\mathcal F$.

Now we let $M=\{f(B(h_1),...,B(h_n))\}$. We have already shown that $\sigma (M)$ is the same as $\mathcal F$. Also, $M$ is closed under pointwise multiplication. Let $E:= \bar{M}$ where is closure is with respect to the $L^2(\Omega)$ norm. Let $H:=E \cap L^\infty(\Omega)$. We have $M \subset H$. It can be shown that $H$ is closed with respect to bounded pointwise convergence and also that $H$ contains all constant functions. Thus, the hypothesis is Dynkin's multiplicative theorem has been satisfied and this implies that $H$ contains all $\sigma (M) = \mathcal F$ measureable functions. But this means that is contains all bounded Borel measureable functions. But by the definition of $H$, this implies that $E$ contains all the bounded borel measureable functions. We know that the all $X\in L^2(\Omega)$ can be approximated in the $L^2(\Omega)$ norm by bounded measureable functions. Also since $E$ is closed and contains the bounded measureable functions, then $E$ itself must be equal to $L^2(\Omega)$. Therefore we conclude that $\bar{M} =E = L^2(\Omega)$.

\end{proof}

\begin{NickQuestion}
	Why is $(D_tF)_{t \ge 0}$ a stochastic processes in $L^2(\Omega ; H)$ where $H=L^2(\mathbb R^+)$? $D_tF$ is defined as
	\[
	D_tF = \sum_{i=1}^n \frac{\partial f}{\partial x}(B(h_1),...,B(h_n))h_i(t)
	\]
	where $h_i \in H$ and $B(h_i) = \int_0^\infty h_i dB_s$ is the standard Wiener Integral. So isnt $D_tF$ real valued since we are evaluating the $h_i(\cdot)$ at $t$? Then shouldn't the $(D_tF)_{t \ge 0}$ be a $L^2(\Omega; \mathbb R)$ valued stochastic process?
	\end{NickQuestion}
\textbf{I believe that this is just a typo in the text}

\begin{NickQuestion}[(Proposition 3.2.4)(All questions have been solved with work below)] $\;$
\begin{enumerate}
\item Why does $D_h(\delta(u))=\langle u, h\rangle + \delta(D_h u)$? (I have answered this below)

\item Also, at the bottom of page 52, why is it that (I have answered this below)
	\[
	E \left( \sum_{i,j=1}^\infty D_{e_i}\langle u, e_j \rangle D_{e_j}\langle v, e_i \rangle  \right)) = E \left( \int_0^\infty\int_0^\infty D_su_tD_tv_s dsdt \right)
	\]

\item In the proof of property 3.3 why do we have  (Answered below)
\[
	E\left(\sum_{i=1}^\infty \langle v,e_i \rangle  \bigg(\langle u, e_i \rangle + \delta(D_{e_i}(u))\bigg) \right) = E(\langle u,v \rangle ) + E \left(\sum_{i,j=1}^\infty D_{e_i}\langle e, e_j \rangle D_{e_j}\langle v,e_i \rangle \right)
\]
\end{enumerate}
\end{NickQuestion}

\textbf{Proof of $D_h(\delta(u))=\langle u, h\rangle + \delta(D_h u).$ }

Since $u=\sum F_jh_j$ we just apply the definition of $\delta$ to get
\begin{align}
D_h(\delta(u)) &= D_h\left(  \sum_{j=1}^n F_j B(h_j) - \sum_{j=1}^n\langle DF_j, h_j \rangle \right) \nonumber
\\&=   \sum_{j=1}^n D_h(F_j B(h_j)) - \sum_{j=1}^nD_h\langle DF_j, h_j \rangle \nonumber
\\&= \sum_{j=1}^n F_j \langle h, h_j \rangle +\sum_{j=1}^n B(h_j)D_h(F_j) - \sum_{j=1}^nD_h\langle DF_j, h_j \rangle  \label{P:3.2.4-1}
\end{align}
where the third inequality follows from the product rule and $D_h(B(h_j)) =  \langle h, h_j \rangle$.
Now we focus on the last term from the last line above \eqref{P:3.2.4-1}.
\begin{align}
\sum_{j=1}^nD_h\langle DF_j, h_j \rangle & = \sum_{j=1}^n\langle D_t \langle DF_j, h_j \rangle, h(t) \rangle \label{P:3.2.4-6}
\end{align}
Before we continue, we simplify $D_t \langle DF_j, h_j \rangle$. We may assume that $F_j=f^j(B(h_1^j(t),...,B(h_n^j(t)))$ where $f^j \in C_p^\infty(\R^n)$ and $h_i^j \in H$.
\begin{align}
D_t \langle DF_j, h_j \rangle &= D_t \langle\sum_{i=1}^n \frac{\partial f^j}{\partial x_i}h_i^j(s), h_j(s) \rangle \nonumber
\\ &= \sum_{i=1}^n\langle h_i^j(s), h_j(s) \rangle D_t\left( \frac{\partial f^j}{\partial x_i}\right) \label{P:3.2.4-2}
\end{align}
We now plug \eqref{P:3.2.4-2} into \eqref{P:3.2.4-6} above to get
\begin{align*}
\sum_{j=1}^n\langle D_t \langle DF_j, h_j \rangle, h(t) \rangle &= \sum_{j=1}^n \sum_{i=1}^n\langle h_i^j(s), h_j(s) \rangle \langle D_t\left( \frac{\partial f^j}{\partial x_i}\right), h(t)\rangle
\\&= \sum_{j=1}^n \sum_{i=1}^n\left\langle h_i^j(s)\langle D_t\left( \frac{\partial f^j}{\partial x_i}\right), h(t)\rangle, h_j(s) \right\rangle
\\&= \sum_{j=1}^n \left\langle \langle D_t\left(\sum_{i=1}^nh_i^j(s) \frac{\partial f^j}{\partial x_i}\right), h(t)\rangle, h_j(s) \right\rangle
\\& = \sum_{j=1}^n \bigg\langle \langle D_t\left(DF_j\right), h(t)\rangle, h_j(s) \bigg\rangle
\\& = \sum_{j=1}^n \bigg\langle D_h\left(DF_j\right), h_j(s) \bigg\rangle
\end{align*}
and if I haven't made a mistake, then what we have shown so far is that
\begin{align}
D_h(\delta(u))  & = \sum_{j=1}^n F_j \langle h, h_j \rangle +\sum_{j=1}^n B(h_j)D_h(F_j) - \sum_{j=1}^nD_h\langle DF_j, h_j \rangle \nonumber
\\&= \sum_{j=1}^n F_j \langle h, h_j \rangle + \sum_{j=1}^n B(h_j)D_h(F_j)- \sum_{j=1}^n \bigg\langle D_h\left(DF_j\right), h_j \bigg\rangle \label{P:3.2.4-3}
\end{align}
Also, the above is exactly line two of the first centered equations on page 53.


Now we need to show that the last two terms add up to $\delta(D_h(u))$. This is not defined in the text but I think that we need to define $D_hu$ as the following
\begin{align}
D_hu = D_h \sum_{j=1}^n F_jh_j(t) =  \sum_{j=1}^n h_j(t) D_h(F_j) \label{P:3.2.4-4}
\end{align}
Thus now $D_h(u)$ \eqref{P:3.2.4-4} is of the form $\sum F_j h_j$ and so we may apply the definition of $\delta$ to it to get the following
\begin{align}
\delta(D_hu) & = \sum_{j=1}^n B(h_j) D_h(F_j) -  \sum_{j=1}^n \langle D_s(D_hF_j), h_j(s) \rangle \label{P:3.2.4-5}
\end{align}
Now we focus in the last term of the right hand side of \eqref{P:3.2.4-5}. From \eqref{P:3.2.4-3} we note that we want to show that
\begin{align}
\sum_{j=1}^n \langle D_s(D_hF_j), h_j(s) \rangle = \sum_{j=1}^n \bigg\langle D_h\left(DF_j\right), h_j \bigg\rangle \label{P:3.2.4-7}
\end{align}
Here are the steps,
\begin{align*}
	\sum_{j=1}^n \langle D_s(D_hF_j), h_j(s) \rangle &= \sum_{j=1}^n \left\langle D_s\left( \left\langle \sum_{i=1}^n \frac{\partial f^j}{\partial x_i} h_i^j(t),h(t) \right\rangle \right), h_j(s) \right\rangle
	\\&  = \sum_{j=1}^n \left\langle D_s\left( \sum_{i=1}^n \frac{\partial f^j}{\partial x_i} \left\langle  h_i^j(t),h(t) \right\rangle \right), h_j(s) \right\rangle
	\\& = \sum_{j=1}^n \left\langle \sum_{i=1}^n D_s\left(  \frac{\partial f^j}{\partial x_i}\right) \left\langle  h_i^j(t),h(t) \right\rangle , h_j(s) \right\rangle
	\\& = \sum_{j=1}^n \left\langle \left\langle \sum_{i=1}^n D_s\left(  \frac{\partial f^j}{\partial x_i}\right)   h_i^j(t),h(t) \right\rangle , h_j(s) \right\rangle
	\\& = \sum_{j=1}^n \left\langle \left\langle \sum_{i=1}^n D_s\left(  \frac{\partial f^j}{\partial x_i}h_i^j(t)\right)  ,h(t) \right\rangle , h_j(s) \right\rangle
	\\& = \sum_{j=1}^n \bigg\langle \big\langle  D_s\left(  D_tF_j\right)  ,h(t) \big\rangle , h_j(s) \bigg\rangle
	\\& = \sum_{j=1}^n \bigg\langle  D_h\left(  DF_j\right) , h_j \bigg\rangle
\end{align*}
Hence \eqref{P:3.2.4-7} is shown and we are done.

\textbf{Proof of bottom of page 52}

\begin{align*}
	\sum_{i,j=1}^\infty D_{e_i}\langle u , e_j \rangle D_{e_j}\langle v, e_i \rangle &= 	\sum_{i,j=1}^\infty \bigg\langle D_s\langle u , e_j \rangle, e_i(s) \bigg\rangle \bigg\langle D_t\langle v, e_i\rangle, e_j(t) \bigg\rangle
	\\& = \sum_{i,j=1}^\infty \bigg\langle \langle D_s(u_t) , e_j(t) \rangle, e_i(s) \bigg\rangle \bigg\langle \langle D_t(v_s), e_i(s)\rangle, e_j(t) \bigg\rangle
		\\& = \sum_{i=1}^\infty \bigg\langle\sum_{j=1}^\infty \left( \langle D_s(u_t) , e_j(t) \rangle \bigg\langle \langle D_t(v_s), e_i(s)\rangle, e_j(t) \bigg\rangle\right) , e_i(s) \bigg\rangle
		\\& = \sum_{i=1}^\infty \bigg\langle  \bigg\langle D_s(u_t) , \big\langle D_t(v_s), e_i(s) \big\rangle \bigg\rangle , e_i(s) \bigg\rangle
		\\&= \sum_{i=1}^\infty \bigg\langle  \bigg\langle D_s(u_t) , \big\langle D_t(v_q), e_i(q) \big\rangle \bigg\rangle , e_i(s) \bigg\rangle
		\\& = \sum_{i=1}^\infty \int_0^\infty \left( \int_0^\infty D_s(u_t) \left( \int_0^\infty D_t(v_q)e_i(q) \ud q\right) \ud t \right) e_i(s) \ud s
		\\& =  \int_0^\infty \left( \int_0^\infty D_s(u_t) \left(\sum_{i=1}^\infty \int_0^\infty D_t(v_q)e_i(q) \ud q \;  e_i(s) \right) \ud t \right)\ud s
		\\& =  \int_0^\infty \int_0^\infty D_s(u_t) D_t(v_s) \ud t \ud s
\end{align*}
Where above we use the identity
\[
	\langle f, g \rangle = \sum_{i=1}^\infty \langle f, e_i \rangle \langle g , e_i \rangle.
\]

\textbf{Proof of step from property 3}


\begin{align}
E\left(\sum_{i=1}^\infty \langle v,e_i \rangle  \bigg(\langle u, e_i \rangle + \delta(D_{e_i}(u))\bigg) \right) &= E(\langle u,v \rangle ) + E\left(\sum_{i=1}^\infty \langle v,e_i \rangle\delta(D_{e_i}(u))\right) \nonumber
\\&= E(\langle u,v \rangle ) + \sum_{i=1}^\infty E\bigg( \bigg\langle D(\langle v,e_i \rangle),D_{e_i}(u)\bigg\rangle\bigg) \label{s-1}
\\&= E(\langle u,v \rangle ) + \sum_{i=1}^\infty E\bigg( \bigg\langle \langle D_s(v_t),e_i(t) \rangle,\langle D_t(u_s),e_i(t) \rangle\bigg\rangle\bigg) \nonumber
\\&= E(\langle u,v \rangle ) + \sum_{i,j=1}^\infty E\left( \bigg\langle \langle D_s(v_t),e_i(t) \rangle,e_j(s)\bigg\rangle \bigg\langle \langle D_t(u_s),e_i(t)\rangle , e_j(s) \bigg\rangle \right) \nonumber
\\&= E(\langle u,v \rangle ) + \sum_{i,j=1}^\infty E\left(  D_{e_j}\big\langle v_t,e_i(t)\big\rangle \bigg\langle \langle D_t(u_s),e_j(s)\rangle , e_i(t) \bigg\rangle \right) \label{s-2}
\\&= E(\langle u,v \rangle ) + \sum_{i,j=1}^\infty E\left(  D_{e_j}\big\langle v_t,e_i(t)\big\rangle D_{e_i}\big\langle u_s, e_j(s)\big\rangle \right) \nonumber
\\&= E(\langle u,v \rangle ) + \sum_{i,j=1}^\infty E\left(  D_{e_j}\big\langle v,e_i\big\rangle D_{e_i}\big\langle u, e_j\big\rangle \right) \nonumber
\end{align}
 where in \eqref{s-1} we used Proposition 3.2.3 and in \eqref{s-2} we use Fubini's Theorem.
\begin{NickQuestion}[(Why is $\delta$ linear?)]
	On page 55 it says that $\delta$ is linear. First, for $u=\sum_{j=1}^nF_jh_j$ we have
	\[
	\delta(u) = \sum_{j=1}F_jB(h_j) - \sum_{j=1}^n \langle DF_j, h_j \rangle_H
	\]
	Now if we have a $v = \sum_{j=1}^nG_jh'_j$ how does $\delta(u+v) = \delta(u) + \delta(v)$. Also how could we write $\sum_{j=1}^nF_jh_j + \sum_{j=1}^nG_jh'_j = \sum_{j=1}^nK_jk_j$ where $K_j \in \mathcal S$ and $k_j \in H$

	Do we do this through the relationship $E(\langle DF,u \rangle ) = E(\delta(u)F)$ in the sense that
	\begin{align*}
	E(\delta(u+v)F) &= E(\langle DF,u +v\rangle )
	\\&=  E(\langle DF,u\rangle ) +  E(\langle DF,v\rangle )
	\\& = E(\delta(u)F) + E(\delta(v)F)
	\\& = E((\delta(u)+\delta(v))F)
	\end{align*}
	and if this is true for every $F$ then I believe we can say  $\delta(u+v) = \delta(u) + \delta(v)$?
	\textcolor{magenta}{Yes, this is the argument to show that $\delta$ is linear.
	In general, for two distributions $F$ and $G$, how one identifies $F=G$? This is true iff $\InPrd{F,\varphi}=\InPrd{G,\varphi}$ for all test function $\varphi$.
	Here, is the same. In order to identify that two random variables $\delta(u+v)$ and $\delta\left(v\right)+\delta(u)$ are same, this relation has to be tested against all functions in its dual
	space, which is $\mathbb{D}^{1,2}$.}

\end{NickQuestion}

\begin{NickQuestion}[Exercise 3.1]
	I don't see how we can apply this hint? Suppose we are able to show that the hint is true then how does that show $\mathcal{P}$ is dense in $L^q(\Omega)$? There is a theorem that $\{x_n\}$ is a
	basis for a Hilbert space $H$ if and only if $\sum \langle x,x_n \rangle x_n = 0$ implies $x=0$ but this is only valid for Hilbert space and $L^q(\Omega)$ for $q \not= 2$ is not a Hilbert
	space since there is no inner product.
	\textcolor{magenta}{This is the standard argument to show one set is dense in another.
		You may miss the fact that  the set of polynomial functions is dense in $C^{\infty}_0(\R^d)$, the latter is dense in the set of the Schwatz test function
		$\mathcal{S}(\R^d)$.
	Please do justify yourself this problem.}
\end{NickQuestion}

\begin{NickQuestion}[Prop 3.3.2 page 54]
	This relates to exercise 3. In the proof of prop 3.3.2 how can we use Holders inequality to show that $\varphi'(F)D(F) \in L^q(\Omega,H)$?
\end{NickQuestion}
\begin{proof}[Nicks Solution]
	\begin{align*}
		\left[E\Norm{\varphi'(F)DF}_H^q\right]^{1/q} &= \left[\int_\Omega \left(\int_0^\infty \varphi'(F)^2(D_tF)^2 \ud t\right)^{q/2}\ud P\right]^{1/q}
		\\& \le \left[\int_\Omega (1+|F|^\alpha)^q \left(\int_0^\infty(D_tF)^2 \ud t\right)^{q/2} \ud P\right]^{1/q}
		\\&= \Norm{(1+|F|^\alpha)\left(\int_0^\infty(D_tF)^2 \ud t\right)^{1/2}}_q
	\end{align*}
	Note that $\dfrac{1}{q} = \dfrac{1}{p} + \dfrac{1}{p/\alpha}$ and so if we apply Holders Inequality above we get that
	\begin{align*}
			\left[E\Norm{\varphi'(F)DF}_H^q\right]^{1/q}  &\le \Norm{1+|F|^\alpha}_{p/\alpha}\Norm{\left(\int_0^\infty(D_tF)^2 \ud t\right)^{1/2}}_p
	\end{align*}
	which is finite since $F \in \mathbb D^{1,p}$.
	\textcolor{magenta}{
		Or equivalently,
		\begin{align*}
			\Norm{\Norm{ \varphi'(F) DF}_H}_q^2
			=& \Norm{\varphi(F)\Norm{DF}_H}_q^2\\
			\le& \Norm{\varphi\left(F\right)}_{\frac{p}{\alpha}}^2 \Norm{\Norm{DF}_T}_p^2 \\
			= & \Norm{\varphi\left(F\right)}_{\frac{p}{\alpha}}^2 \int_0^\infty \Norm{D_t F}_p^2 \ud t
		\end{align*}
	}
\end{proof}

\begin{NickQuestion}[Prop 3.3.2 page 54]
	I have not been able to verify why $\varphi(F) \in \mathbb D^{1,q}$. By definition, this means that we need to show that there exists $X_n \in \mathcal S$ such that
	\begin{enumerate}
		\item $X_n \to \varphi(F)$ in $L^q(\Omega)$
		\item $D(X_n) \to D(\varphi(F))$ in $L^q(\Omega,H)$
	\end{enumerate}

	\textcolor{magenta}{You have made the problem more complex than what it should be. You only need to check the seminorm $\Norm{\cdot}_{1,p}$ is finite.}


	By assumption we already have that there exists $\{F_n\} \subset \mathcal S$ such that $F_n \to F$ in $L^p(\Omega)$ and also $DF_n \to DF$ in $L^p(\Omega,H)$. We then take a sequence of
	mollifiers, $\{\alpha_n\}$ and consider the functions $\varphi_k(x) = (\varphi * \alpha_k)(x)$ which by the original assumption on $\varphi$ are $\alpha+1$ holder continuous.
	\textcolor{magenta}{You cannot say a function is $\beta$-H\"older continuous with $\beta\ge 1$. That means the function is a constant.}

	Now, we notice that by properties of the mollifier since $\varphi(F) \in L^q(\Omega)$ then
		\begin{align*}
			E|\varphi_k(F_j)-\varphi(F_j)|^q \xrightarrow[\text{\text{$j \to \infty$}}]{\text{}} 0
		\end{align*}
	Next, since $\varphi$ is $\alpha +1$ holder continuous then we have
		\begin{align*}
				E|\varphi(F_j)-\varphi(F)|^q &\le KE|F_j-F|^{q(\alpha+1)}
				\\&=KE|F_j-F|^{p} \to 0
		\end{align*}
	and thus we conclude as $k$ and $j$ go to infinity that $\varphi_k(F_j) \to \varphi(F)$ in $L^q(\Omega)$ and $\varphi_k(F_j)\in \mathcal S$.

	I am unable to show why $D_t(\varphi_k(F_j)) \to D_t(\varphi(F))$ in $L^q(\Omega,H)$.

	However since $F_j=f^j(B(h^j_1),...,B(h^j_n))$, we see that
		\begin{align*}
			D_t(\varphi_k(F_j)) &= \sum_{i=1}^n \dfrac{\partial}{\partial x_i}(\varphi * \alpha_k)(F_j)h^j_i(t)
			\\&= \sum_{i=1}^n (\varphi' * \alpha_k)(F_j)\dfrac{\partial f^j}{\partial x_i}h^j_i(t)
			\\& =  (\varphi' * \alpha_k)(F_j)\sum_{i=1}^n \dfrac{\partial f^j}{\partial x_i}h^j_i(t)
			\\& = (\varphi' * \alpha_k)(F_j)D_t(F_j).
		\end{align*}
	We know that $\{\alpha_k\}$ are mollifiers so $ (\varphi' * \alpha_k)(F_j) \to \varphi '(F_j)$ almost surely and by assumption $D_t(F_j) \to D_t(F)$ in $L^p(\Omega,H)$. But I still am having
	trouble rigorously showing that $	D_t(\varphi_k(F_j)) \to 	D_t(\varphi(F_j))$ in $L^p(\Omega, H)$.
	\textcolor{magenta}{You have two indices, you need to take the limit once for just one indices, not the two at the same time. More precisely,
	\begin{align*}
		\Norm{D_t\left(\varphi_k(F_j) - D_t\left(\varphi(F)\right)\right)}_p
		\le \Norm{D_t\left(\varphi_k(F_j)\right)- D_t\left(\varphi(F_j)\right)}_p
		+  \Norm{D_t\left(\varphi(F_j)\right)- D_t\left(\varphi(F)\right)}_p.
	\end{align*}
	Actually, you need to split to three such differences. I will explain to you during the meeting.
	}
\end{NickQuestion}

\begin{NickQuestion}[Exercise 3.4]
	Let $M = \sum_{t\in [0,1]} X_t$. The hint says to approximate this by taking the sup over a finite set. We know that since $X_t$ is continuous in $t$ then $M=\sup_{t\in \mathbb A}X_t$ where
	$A=\mathbb{Q} \cap [0,1]$. Let $\{a_n\}$ be an enumeration of $A\setminus \{0,1\}$ and let $S_k = \{0,a_1,...,a_k,1\}$ then each $S_k$ are finite and $M_k := \sup_{t \in S_k}X_t$ is an
	\textcolor{magenta}{(Note how to type $\mathcal{Q}$, you need curly braces.)}
	approximation of $M$ since $M_k \to M$ almost surely. The hint also tells us to use Corollary 4.2.5 but that requires us to find $D_tM_k$ which I am unsure of how to do. The hint says that we
	need to use exercise 3.3 which is chain rule for Lipschitz functions but I do not know how that applies. This corollary would also require us to show that
	$M_k \stackrel{L^2(\Omega)}{\longrightarrow} M$ which I have been unable to show either.
	\textcolor{magenta}{Try harder... This is an interesting problem that is related the hitting problems that Robert Dalang and Eulalia have done a lot work in the SPDE setting.}
\end{NickQuestion}

\subsection{Exercises}

\textbf{Exercise 3.3} Let $\alpha_n$ be a sequence of mollifiers and $\varphi_k := (\varphi * \alpha_k)$. Since $\varphi$ is Lipschitz with Lipschitz constant $K$ then this implies that $\varphi_k ' \le K$ and in addition $\varphi '$ exists almost everywhere and is also bounded by $K$. By proposition 3.3.2 we have that
\[
	D_t(\varphi_k(F))=\varphi_k'(F)D_t(F)
\]
Also note that
\begin{align*}
	|\varphi_k(F)-\varphi(F)|^2 &= \left| \int_\R (\varphi(F-y)-\varphi(F))\alpha_k(y) \ud y \right|^2
	\\ &\le K \left| \int_\R y \alpha_k(y) \ud y \right|^2
	\\ & \le K'
\end{align*}
Hence 	$|\varphi_k(F)-\varphi(F)|^2$ is bounded above and also converges to $0$ almost surely so by the dominated convergence theorem we have that
\[
	\mathbb E |\varphi_k(F)-\varphi(F)|^2 \to 0
\]
and so the hypothesis of corollary 4.2.5 is satisfied. This implies that $\varphi(F) \in \mathbb D^{1,2}$ and
\[
D_t(\varphi_n(F)) \to D_t(\varphi(F))
\]
 weakly in $L^2(\Omega,H)$. Lastly, since $\varphi_k'(F)D_t(F)$ is bounded in $L^2(\Omega,H)$ and also
 \[
 \varphi_k'(F)D_t(F) \to G D_t(F)
 \]
 almost surely where $G$ is a random variable bounded by $K$ then
 \[
 \varphi_k'(F)D_t(F) \to G D_t(F)
 \]
 weakly in $L^2(\Omega,H)$ as well. Since weak limits are unique then we must have
\[
	D_t(\varphi(F)) = G D_t(F)
\]
However note that in general we have as a property of mollifiers that $\varphi_n'(x) \to \varphi'(x)$ almost surely. But we said before that $\varphi'(x)$ only exists almost surely since $\varphi$ is
Lipschitz. Therefore if the law of $F$ is absolutely continuous with respect to Lebesgue measure then we have $G=\varphi '(F)$ almost surely. Indeed, let $\mathcal N$ be the null set where $\varphi'$
is undefined. If the law of $F$ is absolutely continuous wit respect to Lebesgue measure then $P(F \in \mathcal N)=0$ and we must have $G=\varphi '(F)$ almost surely.

\textbf{Exercise 3.4} I found a variant of this problem $\href{http://www.math.wisc.edu/~kurtz/NualartLectureNotes.pdf}{here}$ where they consider $\{B_t\}_{t\in[0,1]}$ instead of an arbitrary centered Gaussian processes $\{X_t; t\in[0,1]\}$ . They use the idea I was trying, which is we want to consider the function $f(x)=\max_{i\in \{1,...,n\}}x_i$ for $x \in \R^n$. $f$ is a Lipschitz function and if we consider $\{q_n\}$ an enumeration of $\mathbb{Q}\cap [0,1]$ and consider $Y_n=(X_{q_1},..., X_{q_n})$ then $M_n:= \max_{i\in \{1,...,n\}}\{X_{q_1},..., X_{q_n}\} = f(Y_n)$ and $M_n \to M$ almost surely. However in the above linked pdf they claim without proof that that $M_n \to M$ in $L^2(\Omega)$ which I have so far been unable to show. 

I was unaware that there is an extension of the chain rule for Lipschitz functions that is given in the above linked pdf that says we may consider Lipschitz functions on $\R^d$ and random vectors $F=(F_1,...,F_n)$ whose components are in $\mathbb {D}^{1,2}$ and then $f(F)$ will also be in $\mathbb {D}^{1,2}$. We apply this above to see that $M_n:= \max_{i\in \{1,...,n\}}\{X_{q_1},..., X_{q_n}\} = f(Y_n) \in\mathbb {D}^{1,2} $. Next we define the sets 
\begin{align*}
	&A_1 = \{M_n = X_{q_1}\}
	\\ &...
	\\ & A_k = \{M_n \not= X_{q_1}, M_n \not= X_{q_2},...,M_n= X_{q_k}\}, 2 \le k \le n
\end{align*}
and by doing so 
\begin{align*}
DM_n &= D\left(\sum_{k=1}^n1_{A_k}X_{q_k}\right)
\\&= \sum_{k=1}^n1_{A_k}D\left(X_{q_k}\right)
\end{align*}
where the last equality is to follow because Lemma 3.4.2 on page 56 but I need to verify why. Lets suppose for now that $\{X_t\}$ are Brownian motions. Then 
\begin{align} \label{brown}
	D\left(X_{q_k}\right) = 1_{[0,q_k]}(t)
\end{align}
and the above becomes 
\begin{align*}
	DM_n &= \sum_{k=1}^n1_{A_k}D\left(X_{q_k}\right)
	\\ & = \sum_{k=1}^n1_{A_k}1_{[0,q_k]}(t)
\end{align*}
and we have that 
\begin{align*}
	\mathbb{E}\Norm{DM_n}_H^2 &= \mathbb{E}\Norm{\sum_{k=1}^n1_{A_k}1_{[0,q_k]}(t)}_H^2
	\\& = \mathbb{E} \int \left( \sum_{k=1}^n1_{A_k}1_{[0,q_k]}(t) \ud t \right)^2
	\\& = \sum_{k=1}^n  \mathbb{E}\left(1_{A_k}\int1_{[0,q_k]}(t) \ud t\right)
\end{align*}
where the last inequality follows from the disjointness of the $A_k$. Continuing the calculations we see 
\begin{align*}
\mathbb{E}\Norm{DM_n}_H^2 &= \sum_{k=1}^n  \mathbb{E}\left(1_{A_k}\int1_{[0,q_k]}(t) \ud t\right)
\\& = \sum_{k=1}^n  q_k \mathbb{E}\left(1_{A_k}\right)
\\& \le \sum_{k=1}^n  \mathbb{E}\left(1_{A_k}\right)
\\& \le \mathbb{E}\left(\Omega \right) =1
\end{align*}
and so we apply Corollary 4.2.5 to conclude that $M=\sup_{t \in [0,1]}X_t \in \mathbb{D}^{1,2}$ and that $DM_n \to DM$ in the weak topology of $L^2(\Omega,H)$ which means that for all $G\in \mathbb{D}^{1,2}$
\[
\langle DM_n,G \rangle_{L^2(\Omega,H)} \to \langle DM,G\rangle_{L^2(\Omega,H)}
\]
To prove that $DX=1_{[0,T]}$ we note that $\sum_{k=1}^n1_{A_k}q_k$ which converges almost surely to $T$. Then we can conclude that $DM_n  = \sum_{k=1}^n1_{A_k}1_{[0,q_k]}(t)$ converges to $1_{[0,T]}$ almost surely. Since $DM_n$ converges both weakly and alost surely then a result from functional analysis says that these limits must be equal and we see that $DM=1_{[0,T]}$. 

Now if we go back to \eqref{brown} and assume that $\{X_t\}$ is a centered Gaussian process then I am unsure of how to continue. For example we can not calculate $D(X_t)$ for any $t\in[0,1]$ unless we know that $X_t = \int_0^\infty h(s)dB_s$ for some $h \in H$ but I do not think that this can be claimed. 


\section{Chapter 4}
\subsection{Reading material} 
\begin{NickQuestion}[proof of proposition 4.2.1]
	In the conclusion of this proof, Naulart refers to a density argument but what is he referring too? We see that from (4.1) on page 64 that 
	\[
		\mathbb{E}I_n^2(f) = n!\Norm{\tilde{f}}_{L^2(\R_+^n)}
	\]
so is there a result that says that fuctions  
\[
A:=\{g^{\otimes n}(x_1,...,x_n) = g(x_1)\cdots g(x_n), g\in L^2(\R^+)\}
\]
 are dense in $L^2_s(\R^n_+)$? If thats the case then I believe that for a given $f \in L^2_s(\R^n_+)$ then we find $\{g_k\}\in A$ such that 
 \[
 	\mathbb{E}|I_n(f-g_k)|^2 = n!\Norm{\tilde{f}-g_k}_{L^2(\R_+^n)} \to 0
 \]
 and I think since $D$ is a closed operator then we need 
 \[
 D_t(I_n(g_k)) = nI_{n-1}(g_k(\cdot,t)) \to D_t(I_n(f))?
 \]
\end{NickQuestion}

\begin{NickQuestion}[Example 4.2.3]
	We know that $B_1^3=\sum_{n=0}^\infty I_n(f_n)$ and by the bottom of page 66, for this example we se use the formula 
	\[
	f_n=\dfrac{1}{n!}\mathbb{E}(D^nB_1^3)
	\]
and calculating each of the above we see that 
\begin{align*}
&f_0 = \mathbb{E}B_1^3=0
\\& f_1 = 31_{[0,1]}(t_1)
\\& f_2 =0
\\& f_3 = 1_{[0,1]}(t_1) 1_{[0,1]}(t_2) 1_{[0,1]}(t_3)
\\& f_k = 0,\; k \ge 4
\end{align*}
thus we see that by using the definition on page 63 that
\begin{align*}
B_1^3 &= I_1(31_{[0,1]}(t_1)) + I_3(1_{[0,1]}(t_1) 1_{[0,1]}(t_2) 1_{[0,1]}(t_3))
\\&= \int_0^\infty 31_{[0,1]}(t_1)) \ud B_{t_1} + \int_0^\infty \int_0^{t_3}\int_0^{t_2}1_{[0,1]}(t_1) 1_{[0,1]}(t_2) 1_{[0,1]}(t_3) \ud B_{t_1} \ud B_{t_2} \ud B_{t_3}
\\&= 3B_1 + \int_0^1 \int_0^{t_3} 1_{[0,1]}(t_2) \int_0^{t_2}1_{[0,1]}(t_1)  \ud B_{t_1} \ud B_{t_2} \ud B_{t_3}
\\&= 3B_1 + \int_0^1 \int_0^{t_3 \wedge 1} \int_0^{t_2 \wedge 1} \ud B_{t_1} \ud B_{t_2} \ud B_{t_3}
\\& =  3B_1 + \int_0^1 \int_0^{t_3} \int_0^{t_2} \ud B_{t_1} \ud B_{t_2} \ud B_{t_3}
\end{align*}
This is different from the answer in the test on page 67. Is mine correct? 
\end{NickQuestion}

\begin{NickQuestion}[directional derivative page 69]
	\begin{enumerate}
\item $\tau_g(\omega) = \omega +g$ where $g$ is continuous on $\R^+$ and $g(0)=0$. Why does this imply that 
\[
	\tau_g(B(h))=B(h) + \langle h, \dot{g} \rangle_H. 
\]
when by definition 
\[
	\tau_g(B(h))=B(h) + g
\]
but also $B(h) \not \in \Omega$ where $\Omega$ is the Weiner space and 
\[
B(h)=\int_0^\infty h(s)\ud B_s
\] 
so why does $\tau_g(B(h))$ exist?

\item Why does Girsonovs theorem imply that if $F \in L^2(\Omega)$ then $F(\omega +g)$ is almost surely defined? 

\item What is the motivation behind the definition of the directional derivative 
\[ 
	\tilde{D}_gF = \lim_{\epsilon \to 0} \dfrac{1}{\epsilon}(F \circ \tau_{\epsilon g} -F)?
\]
$\href{http://www.math.wisc.edu/~kurtz/NualartLectureNotes.pdf}{\text{In this pdf}}$ they claim that when we have $DF = \sum_{i=1}^n\partial_i f(W(h_1),...,W(h_n))h_i$ then the directional derivative is 
\begin{align} \label{pdf1}
\langle DF, h \rangle_H = \lim_{\epsilon \to 0}\frac{1}{\epsilon}(f(W(h_1)+\epsilon\langle h_1,h \rangle,...,W(h_n)-\epsilon \langle h_n,h \rangle)-F)
\end{align}
and also it clams that $\langle DF, h \rangle_H $ is the derivative at $\epsilon=0$ of the shifted processes 
\begin{align} \label{pdf2}
F(W(g)+\epsilon \langle g,h\rangle )
\end{align}
for $g\in H$ but what does this mean?

Consider $F=(W(h_1),...,W(h_n))$ and here we have 
\[
DF =\sum_{i=1}^n \frac{\partial f}{\partial x_i}(W(h_1),...,W(h_n))h_i
\]
and so if we try to calculate $\langle DF ,h \rangle$ we get 
\begin{align} \label{ddiv}
	\langle DF ,h \rangle &=  \sum_{i=1}^n\frac{\partial f}{\partial x_i}(W(h_1),...,W(h_i),...,W(h_n))\langle h_i, h \rangle   
	\\&= \sum_{i=1}^n\lim_{\epsilon \to 0} \dfrac{f(W(h_1),...,W(h_i)+\epsilon \langle h_i ,h \rangle, ,...,W(h_n))-F}{\epsilon \langle h_i, h \rangle } \langle h_i , h \rangle \nonumber
	\\& = \sum_{i=1}^n df\lim_{\epsilon \to 0} \dfrac{f(W(h_1),...,W(h_i)+\epsilon \langle h_i ,h \rangle, ,...,W(h_n))-F}{\epsilon } \nonumber
\end{align}
but I do not see how we can simplify this to \eqref{pdf1}? 

	\end{enumerate}
\end{NickQuestion}

\textbf{Solution for 3 \eqref{pdf1}:} The definition of the directional derivative is 
\[
	D_v(f)=\lim_{h \to 0} \dfrac{f(x+hv)-f(x)}{h}  = \dfrac{d}{d\epsilon} f(x+\epsilon v) |_{\epsilon=0}
\]  
and according to wikipedia, when the partials of $f$ exist then it can be shown that 
\begin{align}
	D_v(f) = \sum_{i=1}^nv_i\dfrac{\partial f}{\partial x_i} \label{dirdiv}
\end{align}
where $v=(v_1,...,v_n)$. This is stated with out proof and is supposed to follow from the definition of the total derivative which is something I am not familiar too. However, if we let $v = (\langle h_1 ,h \rangle, ..., \langle h_n,h \rangle)$ then 
\begin{align*}
		D_v(f) &=  \lim_{\epsilon \to 0}\frac{1}{\epsilon}(f(W(h_1)+\epsilon\langle h_1,h \rangle,...,W(h_n)-\epsilon \langle h_n,h \rangle)-F)
	\\& = \sum_{i=1}^n \langle h_i,h \rangle \dfrac{\partial f}{\partial x_i}
	\\ & = \langle DF,h \rangle 
	\end{align*}
	where the second equality follows from \eqref{dirdiv} and the third equality follows from the work listed in the gray text from the question I asked. 
\begin{NickQuestion}[Proof of Theorem 4.4.2]
\begin{enumerate}
\item	Why does 
	\[
		\mathbb{E}\left(\left|\frac{1}{\epsilon}(F_n\circ \tau_{\epsilon g}-F_n)-\langle DF_n,\dot{g} \rangle_H \right|\right) \to 0
	\]
for fixed $n$ and $\epsilon \to 0$? 

I believe that they are claiming without proof that if $F_n \in \mathcal S$ then 
\[
\dfrac{1}{\epsilon}(F_n\circ \tau_{\epsilon g} - F_n) \to \langle DF_n,\dot{g} \rangle_H
\]
 I assume in $L^2(\Omega)$ as $\epsilon \to 0$. 
 
 I believe this is true in the most simple example where $F_n=\int_0^Th_n(s)\ud B_s$ where $h_n \in H$ and $B_s$ is the canonical brownian motion on the Wiener space, ie $B_s(\omega)=\omega(s)$. For any $g \in \mathcal H$, or in other words 
 \[
 g(t) = \int_0^t \dot{g}(s)\ud s
 \]
then we can easily calculate the following derivative 
\begin{align*}
	\tilde{D}_{g}(F_n)(\omega) = \lim_{\epsilon \to 0} \frac{1}{\epsilon}(F_n \circ \tau_{\epsilon g}-F_n) &=  \lim_{\epsilon \to 0} \frac{1}{\epsilon}\left( \int_0^t h_n(s) \ud B_s(\omega + \epsilon g) - \int_0^t h_n(s) \ud B_s(\omega) \right) 
	\\&= \lim_{\epsilon \to 0} \frac{1}{\epsilon}\left( \int_0^t h_n(s) \ud (\omega(s) + \epsilon g(s)) - \int_0^t h_n(s) \ud (\omega(s)) \right)
	\\&=   \lim_{\epsilon \to 0} \left( \int_0^t h_n(s) \ud (g(s)) \right)
	\\& = \int_0^t h_n(s) \dot{g}(s)\ud s 
\end{align*}
and from this we conclude that 
\[
		\tilde{D}_{g}(F_n)(\omega) = \langle DF_n , \dot{g} \rangle_H.
\]
	Can we extend this to all $F_n \in \mathcal S$? 
\item Why does 
\[
	\dfrac{1}{\epsilon}(F_n\circ \tau_{\epsilon g}-F_n)=\frac{1}{\epsilon}\int_0^\epsilon \langle DF_n \circ \tau_{\alpha g}, \dot{g} \rangle_H \ud \alpha 
\]
when $F_n \in \mathcal S$. Above I showed that for $F=\int_0^Th(s) \ud B_s$ that we have 
\[
		\dfrac{1}{\epsilon}(F_n\circ \tau_{\epsilon g}-F_n)= \langle h, \dot{g} \rangle_H
\]
Also since 
\[
	DF \circ \tau_{\epsilon g} = h(t)\circ \tau_{\epsilon g} = h(t)
\]
because in this case $DF$ is constant in $\omega$ then we see that 
\[
\langle h, \dot{g} \rangle_H = \dfrac{1}{\epsilon}	\int_0^\epsilon \langle DF_n \circ \tau_{\alpha g}, \dot{g} \rangle_H \ud \alpha
\]
so the statment is true for the most simple $F \in \mathcal S$. However I do not know how to show this for arbitrary $F \in \mathcal S$. 

\item I don't see how Girsanov's theorem is used to conclude the convergence on the top of page 71? Also how is Girsanov's theorem used in steps 3 and 4 on page 72?
\end{enumerate}
\end{NickQuestion}
\textbf{solution to 1:}

The question is whether or not $F\in \mathcal S$ implies that the directional derivative exists and whether or not
\[
	\tilde{D}_g(F) = \langle DF, \dot{g} \rangle_H,\quad \text{where }  g=\int \dot{g}
\]
I showed above in the gray text that this is true when $F$ takes on the most simply form of 
\[
	F = \int_0^T h(s) \ud B_s
\]
where $h \in H$ and $B_s$ is the canonical brownian motion on the Weiner space. However we want this to be true for any general $F\in \mathcal S$ or in other words $F=f(B(h_1),...,B(h_n))$ where $f \in C^\infty_p(\R^n)$. However, such $f$ can be uniformly approximated by functions of the form $x^n$ so I believe that if we can show that for smooth random variables
\[
F=	\left(\int_0^T h(s) \ud Bs\right)^n 
\]
that 
\[
	\tilde{D}_g(F) = \langle DF, \dot{g}\rangle
\]
then we are done. However this is easily shown as follows 
\begin{align*}
	\dfrac{F \circ \tau_{\epsilon g}-F }{\epsilon} &= \dfrac{\left(\int_0^T h(s)\ud B_s(\omega +\epsilon g \right)^n-\left(\int_0^T h(s) \ud B_s(\omega) \right)^n }{\epsilon}
	\\&= \dfrac{\left(\int_0^T h(s)\ud(\omega(s)) + \epsilon \int_0^T h(s) \ud(g(s))\right)^n-\left(\int_0^T h(s) \ud (\omega(s)) \right)^n }{\epsilon}
	\\& = \dfrac{\sum_{i=1}^n {n\choose i }\epsilon^{i} \left(\int_0^T h(s)\ud(\omega(s)) \right) ^{n-i}\left(\int_0^T h(s) \dot{g}(s) \ud s \right)^{i}}{\epsilon}
	\\&= n\left(\int_0^T h(s) \ud(\omega(s))\right)^{n-1}\langle h, \dot{g} \rangle_H +  \sum_{i=2}^n {n\choose i }\epsilon^{i-1} \left(\int_0^T h(s)\ud(\omega(s)) \right) ^{n-i}\left(\int_0^T h(s) \dot{g}(s) \ud s \right)^{i}
	\\& \to n\left(\int_0^T h(s) \ud(\omega(s))\right)^{n-1}\langle h, \dot{g} \rangle_H \quad \text{as } \epsilon \to 0
	\\& = \langle DF , \dot{g} \rangle_H
\end{align*}
which is what we wanted to show. 

\textbf{Partial solution for 2:} Suppose that $F=\left(\int_0^\infty h(t)\ud B_t\right)^2$. Then we can calculate the following
\begin{align}
	F\circ \tau_{\epsilon g} - F &= \left(\int_0^\infty h(t) \ud B_t(\omega + \epsilon g)\right)^2 -  \left(\int_0^\infty h(t) \ud B_t(\omega )\right)^2 \nonumber
	\\&= 2\epsilon \int_0^\infty h(t) \ud B_t(\omega)\langle h, \dot{g} \rangle + \epsilon^2 \langle h, \dot{g} \rangle^2 \nonumber
	\\&= \epsilon\langle DF, \dot{g} \rangle + \epsilon^2 \langle h, \dot{g} \rangle^2\label{lhs}
\end{align}
In addition, we can also calculate 
\begin{align}
\int_0^\epsilon \langle DF\circ\tau_{\epsilon g} ,\dot{g}\rangle \ud \alpha &= \int_0^\epsilon \int_0^\infty D_sF\circ \tau_{\epsilon g} \dot{g}(s) \ud s \ud \alpha \nonumber
\\& =  \int_0^\epsilon \int_0^\infty \left(2h(s)\int_0^\infty h(t) \ud B_t(\cdot)\right)\circ \tau_{\epsilon g} \dot{g}(s) \ud s \ud \alpha \nonumber
\\&= \int_0^\epsilon \int_0^\infty \left(2\int_0^\infty h(t) \ud B_t(\omega + \epsilon g)\right) h(s)\dot{g}(s) \ud s \ud \alpha \nonumber 
\\&= \int_0^\epsilon \int_0^\infty \left(2\int_0^\infty h(t) \ud B_t(\omega) + 2\epsilon\int_0^\infty h(t)\dot{g}(t) \ud t\right)h(s) \dot{g}(s) \ud s \ud \alpha \nonumber 
\\&= \int_0^\epsilon \int_0^\infty 2\int_0^\infty h(t) \ud B_t(\omega)h(s) \dot{g}(s) \ud s \ud \alpha  + \int_0^\epsilon \int_0^\infty2\epsilon\langle h,\dot{g} \rangle h(s) \dot{g}(s) \ud s \ud \alpha \nonumber 
\\&=\epsilon \langle DF,\dot{g} \rangle + \epsilon^2\langle h, \dot{g}\rangle^2 \label{rhs} 
\end{align}
thus we see that $\eqref{rhs}$ and $\eqref{lhs}$ imply that 
\[
	\frac{1}{\epsilon}	\int_0^\epsilon \langle DF\circ\tau_{\epsilon g} ,\dot{g}\rangle \ud \alpha = \frac{1}{\epsilon}(F\circ \tau_{\epsilon g} - F )
\]
and so the claim is true for $F= \left(\int_0^\infty h(t) \ud B_t\right)^2$ which leads us to believe that it is true for $F= \left(\int_0^\infty h(t) \ud B_t\right)^n$ which again leads us to believe that it is true for all $F\in \mathcal S$. 

\textbf{Informal explanation for some of Step 4:}
The only line that I can't quite convince myself is the first centered equation with says 
\[
	\frac{1}{\alpha} \left(E((F\circ \tau_{\alpha g})G)-E(FG)\right) = E\left(F \frac{1}{\alpha}\left((G\circ \tau_{-\alpha g})\exp(\alpha B(h)-\frac{1}{2}\alpha^2 \Norm{h}_H^2)-G\right)\right)
\]
I dont understand why we pass the $\tau_{\alpha g}$ from $F$ to $G$ and it then becomes a $\tau_{-\alpha g}$. Also is not formally stated but the first expectation is under $Q$ and the second is under $P$ with $Q(A)=E(1_AL_t)$ and here 
\[
	L_t = \exp\left( \int_0^t \dot{g_s} \ud B_s - \frac{1}{2}\int_0^t \dot{g}_s^2 \ud s \right)
\]
where $\dot{g} = h$ which is stated on the previous page. 

The next three centered equations are 
\begin{align}
&\lim_{\alpha \to 0} \frac{1}{\alpha}\left((G\circ \tau_{-\alpha g})\exp(\alpha B(h)-\frac{1}{2}\alpha^2 \Norm{h}_H^2)-G\right) \nonumber
\\&= \frac{d}{d\alpha} \left( (G\circ \tau_{-\alpha g})\exp(\alpha B(h)-\frac{1}{2}\alpha^2 \Norm{h}_H^2)) \right) \bigg\vert_{\alpha = 0}
\\&= B(h)G - \langle DG,h \rangle_H 
\\&= \delta(Gh)
\end{align}
The first equality follows from the limit definition of the derivative. The second equality follows from products rule. We have 
\begin{align*}
	 &\frac{d}{d\alpha} \left( (G\circ \tau_{-\alpha g})\exp(\alpha B(h)-\frac{1}{2}\alpha^2 \Norm{h}_H^2)-G) \right) \bigg\vert_{\alpha = 0}
	 \\&= \left( \frac{d}{d\alpha}  (G\circ \tau_{-\alpha g})\right)\exp(\alpha B(h)-\frac{1}{2}\alpha^2 \Norm{h}_H^2)) \bigg\vert_{\alpha = 0} +  (G\circ \tau_{-\alpha g})\left(  \frac{d}{d\alpha}  \exp(\alpha B(h)-\frac{1}{2}\alpha^2 \Norm{h}_H^2)\right) \bigg\vert_{\alpha = 0}
\end{align*}
Note that since $G \in \mathcal S$ then we have that $G=g(B(h_1),...,B(h_n))$ and therefore 
\[
G\circ \tau_{-\alpha g} = g(B(h_1),...,B(h_n))(\omega - \alpha g) = g(B(h_1)-\alpha \langle h_1, \dot{g} \rangle,...,B(h_n)-\alpha \langle h_n, \dot{g}\rangle)
\]
and from this we can calculate $d/d\alpha$
\begin{align*}
	  \frac{d}{d\alpha}  (G\circ \tau_{-\alpha g}) \bigg \vert_{\alpha = 0} &= \lim_{k \to 0} \frac{1}{k}\left(g(B(h_1)-(\alpha +k) \langle h_1, \dot{g} \rangle,...,B(h_n)-(\alpha+k) \langle h_n, \dot{g}\rangle) -G\right)  \bigg \vert_{\alpha = 0}
	  \\&= -\langle DG, \dot{g} \rangle 
\end{align*}
where the last inequality follows from the work shown in \eqref{ddiv} which I believe can e applied because again $G \in \mathcal S$. Lastly, we simply take the derivative to see
\begin{align*}
	\left(  \frac{d}{d\alpha}  \exp(\alpha B(h)-\frac{1}{2}\alpha^2 \Norm{h}_H^2)\right) \bigg\vert_{\alpha = 0} = B(h)
\end{align*}
putting this together give us 
\begin{align*}
		 &\frac{d}{d\alpha} \left( (G\circ \tau_{-\alpha g})\exp(\alpha B(h)-\frac{1}{2}\alpha^2 \Norm{h}_H^2)-G) \right) \bigg\vert_{\alpha = 0} &= B(h)G - \langle DF, \dot{g} \rangle_H
		 \\& = B(h)G - \langle DF, h \rangle_H
\end{align*} 
which completes the verification for step 4. 
\begin{NickQuestion}[The formula  $$I_n(f)I_m(g) = \sum_{r=0}^{n \wedge m}r! {n \choose r}{m \choose r}I_{n+m -2r}(f \otimes_r g)$$ on page 64:] $\;$
	
	This is not a question, just wanted to test this formula out for the most basic example, for $f,g \in L^2(\R_+)$ and $I_1(f)I_1(g)$ so I could have an example saved. 
	
	Using the formula we would need to get 
	\begin{align}
	I_1(f)I_1(g) &= I_2(f\otimes_0 g)+ I_0(f \otimes_1 g) \label{eqn}
	\\&= 2\int_0^\infty \int_0^{t_2}f(t_1)g(t_2)\ud B_{t_1}\ud B_{t_2} +  \int_0^\infty f(x)g(x) \ud x
	\end{align}
	where $(f \otimes_r g)$ is defined on page 64. 
	
	But suppose that we do the following:
	\begin{align*}
		X_t &= \int_0^t f(s) \ud B_s
		\\ Y_t &= \int _0^t g(s) \ud B_s
	\end{align*}
	then by applying Ito's formula we see that 
	\begin{align*}
	X_tY_t  &= \int_0^t X_s \ud Y_s + \int_0^t Y_s \ud X_s + \int_0^t \ud X_s \ud Y_s 
	\\& = \int_0^t \left(\int_0^s f(\xi) \ud B_{\xi}\right)g(s) \ud B_s + \int_0^t \left(\int_0^s g(\xi) \ud B_{\xi}\right)f(s) \ud B_s + \int_0^t f(s)g(s) \ud s
	\\& = \int_0^t \left(\int_0^s f(\xi)g(s)\ud B_{\xi}\right) \ud B_s + \int_0^t \left(\int_0^s g(\xi)f(s) \ud B_{\xi}\right) \ud B_s + \int_0^t f(s)g(s) \ud s
	\\& = \int_0^t \int_0^s f(\xi)g(s)+g(\xi)f(s) \ud B_{\xi} \ud B_s  + \int_0^t f(s)g(s) \ud s
	\\& = \int_0^t \int_0^s 2(f \tilde{\otimes}_0 g)(\xi,s)\ud B_{\xi} \ud B_s  + \int_0^t f(s)g(s) \ud s
	\end{align*}
	where above $(f \tilde{\otimes}_0 g)$ is the symmetrization of $(f \otimes_0 g)$. Thus by letting $t \to \infty$ we see that 
	\begin{align}
	\lim_{t\to \infty} 	X_tY_t = \int_0^\infty f(s) \ud B_s \int_0^\infty g(s) \ud B_s &= I_1(f)I_(g) \nonumber
	\\& = \int_0^\infty \int_0^s 2(f \tilde{\otimes}_0 g)(\xi,s)\ud B_{\xi} \ud B_s  + \int_0^\infty f(s)g(s) \ud s \nonumber
	\\&  = 2I_1(f \tilde{\otimes}_0 g) + I_0(f \otimes_1) \nonumber
	\\& = 2I_1(f \otimes g) + I_0(f \otimes_1) \label {cont}
	\end{align}
	where the last line follows that $I_n(f) = I_n(\tilde{f})$ for any $f \in L^2(\R_+^n)$. 
\end{NickQuestion}

\textbf{Exercises:}

\textbf{Exercise 2:} We want to prove the following for $f,g \in L^2_s(\R_+^2)$
\[
I_2(f)I_2(g) =	I_4(f\otimes_0 g) + 4I_2(f \otimes_1 g) + 2 I_0(f \otimes_2 g).
\]
The hint tells us to use Ito's formula. Consider the following, 
\begin{align*}
	X_t &= \int_0^t \int_0^\xi f(z, \xi) \ud B_z \ud B_{\xi} = \int_0^t F(\xi) \ud B_\xi
	\\Y_t &=  \int_0^t \int_0^\xi g(z, \xi) \ud B_z \ud B_{\xi} = \int_0^t G(\xi) \ud B_\xi
\end{align*}
where we let 
\[
F(\xi)= \int_0^\xi f(z, \xi) \ud B_z \; \quad \; G(\xi)= \int_0^\xi g(z, \xi) \ud B_z
\]
and so 
\[
\lim_{t \to \infty} X_tY_t = I_2(f)I_2(g).
\]
If we apply Ito's formula, we see that 
\begin{align}
	X_tY_t = \int_0^t X_s \ud Y_s + \int_0^t Y_s \ud X_s + \int_0^t \ud X_s \ud Y_s := L_1 + L_2 + L_3 \label{ito1}
\end{align}
If this is correct, then what does $\ud X_s$ and $\ud Y_s$ equal? For, example, is $ \ud X_s = F(s) \ud B_s = \left(\int_0^s f(z,s) \ud B_z\right) \ud B_s$? If this is right, then we can continue with \eqref{ito1}. 
$ \; \\$ 

\noindent \underline{We will just look at $L_1$ right now.}
\begin{align*}
L_1 &=  \int_0^t X_s \ud Y_s 
\\&=  \int_0^t\underbrace{\left( \int_0^s F(\xi) \ud B_\xi \right)}_{X_s} G(s) \ud B_s
\end{align*}
notice the integrand is a product of two processes. The only way I think I can continue is to apply Ito's formula again. But the problem gets messy here. 
\begin{align*}
	X_sG_s = \int_0^s X_k\ud G_k + \int_0^s G_k \ud X_k+ \int_0^s \ud X_k \ud G_k
\end{align*}
To continue, we need to calculate $dG_k$ which can be done with the following Stochastic Leibniz Rule Theorem which is found $\href{https://math.stackexchange.com/questions/2677654/stochastic-leibniz-rule}{here}$. However I believe we need to assume that $g$ is differentiable in addition to being square integrable. According to Leibniz rule we have that 
\[
\ud G_k = \left( \int_0^k\frac{\partial g}{\partial k} (z,k) \ud B_z \right) \ud k + g(k,k)\ud B_k	
\]
I dont think continuing with this strategy gets us to the desired answer though.

\section{Chapter 5:}
\begin{NickQuestion}[Proposition 5.1.2]
	\begin{enumerate}	
		\item What does $F(e^{-t}B+\sqrt{1-e^{-2t}}B')$ mean? What are $B$ and $B'$? Why is he dropping the subscripts $B_t$ and $B'_t$? I thought this was only notation but it seems that something else is happening as well? It seems like $F$ is being evaluated at $e^{-t}B+\sqrt{1-e^{-2t}}B'$ but this can't be possible since $e^{-t}B+\sqrt{1-e^{-2t}}B' \in \R$ and $F \in L^2(\Omega)$. David Naulart $\href{https://bgsmath.cat/wp-content/uploads/2019/02/Course.pdf}{\text{has this pdf}}$. Page 5 and page 20 have some info. 
		\item How does \[ E'(|F(e^{-t}B+\sqrt{1-e^{-2t}}B')|^p)  = E|F|^p ?\]  
		\item Where does \[ E'\left(\exp\left(e^{-t}\lambda B(h)+\sqrt{1-e^{-2t}}\lambda B'(h) -\frac{1}{2}\lambda^2\right)\right) \] come from? 
	\end{enumerate}
\end{NickQuestion}

\begin{NickQuestion}
	Why does Mehler's formula imply that the operator $T_t$ is nonnegative? 
\end{NickQuestion}

\begin{NickQuestion}[Theorem 5.1.3]
	There is a $\href{http://www.numdam.org/article/AIHPB_1976__12_2_105_0.pdf}{\text{refrence here}}$ however it is in French. 
	\begin{enumerate}
		\item Why is $|T_tF| \le |F|$? All I think that we can say is that \[|T_tF| \le \sum_ne^{-nt}|I_n(f_n)| \le \sum_n |I_n(f_n)|\] and \[ |F| \le \sum_n|I_n(f_n)| \] but this does not mean that \[ |T_t(F)| \le |F|. \]
		\item Why does $|T_tF| \le |F|$ imply that we my assume that $F$ and $G$ are nonnegative? 
		\item I think there is some typo in the claim that \[ E(\tilde{E}(F(\hat{B}))G) = E(F(\hat{B})G) \]
		This is because $E(\tilde{E}(F(\hat{B}))G) \in \R$ while $E(F(\hat{B})G) \in \R^{\Omega}$. I am not sure what the correction should be. 
	\end{enumerate}
\end{NickQuestion}

\begin{NickQuestion}[Lemma 5.14]
	\begin{enumerate}
		\item The first line uses Stroock's formula however this can only be used if $F \in \mathbb{D}^{\infty,2}$ and in this lemma we only assume that  $F \in \mathbb{D}^{k,2}$.
		
	\end{enumerate}
\end{NickQuestion}

\begin{NickQuestion}[Theorem 5.3.1]
	\begin{enumerate}
		\item On the top of page 81 how do we get 
		\[
			E(\langle u, DG, \rangle_H) = - \int_0^\infty E(\langle LT_tu,DG \rangle_H) \ud t ?
		\]
		
			\item How can we show this 
		\[ 
		\tilde{E}\left( (D^2G)(e^{-t}B + \sqrt{1-e^{-2t}}\tilde{B}) \right) = \frac{1}{1-e^{-2t}}\tilde{E}\left( \tilde{D}^2(G(e^{-t}B+\sqrt{1-e^{-2t}}\tilde{B})) \right)
		\]
	This is how I am interpreting this. I am guessing that $G(e^{-t}B + \sqrt{1-e^{-2t}}\tilde{B})$ has the following wiener chaos expansion 
	\begin{align*}
		G = \sum_{n=0}^\infty I_n(f_n)
	\end{align*} 
	where we assume that $\hat{B} = e^{-t}B + \sqrt{1-e^{-2t}}\tilde{B}$ which means 
	\begin{align*}
		I_n(f_n) &= \int_0^\infty \cdots \int_0^{z_2} f(z_1,...,z_n)\ud\hat{B}_{z_1}...\ud\hat{B}_{z_n}
		\\& = \int_0^\infty \cdots \int_0^{z_2} f(z_1,...,z_n)\ud \left(e^{-t}B_{z_1} + \sqrt{1-e^{-2t}}\tilde{B}_{z_1}\right)...\ud \left(e^{-t}B_{z_n} + \sqrt{1-e^{-2t}}\tilde{B}_{z_n}\right)
	\end{align*}
	but something does not seem right about this.
	\end{enumerate}
\end{NickQuestion}
\textbf{Solution for 1}
		So we have that $u \in \mathbb{D}^{1,p}(H)$ and so $u=u_s(\omega)$. I believe we can say that $u$ has Weiner Expansion
		\[
			u_s = \sum_{n=0}^\infty I_n(f_{n,s})
		\]
		with $f_{n,s} \in L^2_s(\R^n)$. Therefore we have 
		\[
			T_t(u_s) = \sum_{n=0}^\infty e^{-nt}I_n(f_{n,s})
		\]
		and in return we then have using the formula for $L$ on page 78
		\begin{align*}
			LT_t(u_s) &= \sum_{n=1}^\infty -ne^{-nt}I_n(f_{n,s})
			\\ &= \sum_{n=1}^\infty \frac{\ud}{\ud t}e^{-nt}I_n(f_{n,s})
		\end{align*}
	Now we integrate to see 
	\begin{align*}
\int_0^\infty LT_t(u_s)\ud t &= \int_0^\infty  \sum_{n=1}^\infty\frac{\ud}{\ud t}e^{-nt}I_n(f_{n,s}) \ud t
\\& =  \sum_{n=1}^\infty I_n(f_{n,s}) [ e^{-nt} ]_0^\infty
\\& = - \sum_{n=1}^\infty I_n(f_{n,s})  = -u_s
\end{align*}

\begin{NickQuestion}[Theorem 5.5.1]
	Why are we able to assume that 
	\[
		\varphi(0) := \int_0^\infty x p(x) \ud x = \frac{1}{2}E|F|?
	\]
	This assumes that $\rho$ is a symmetric across the y-axis  but this was never assumed to be the case? 
\end{NickQuestion}

\section{Chapter 6}

\begin{NickQuestion}[page 88]
	How can we show 
	\[
		E\left(1_{\{\sup_{t \le s \le 1}B_s-B_t \ge \sup_{0 \le s \le t}B_s-B_t\}} | \mathcal F_t\right) = 2-2\Phi\left( \frac{\sup_{0 \le s \le t}(B_s-B_t)}{\sqrt{1-t}} \right)
	\]
	
The reflection principle states that for a Brownian motion, $B_t$, and if we define $M_t=\sup_{0 \le s \le t}B_s$ then 
\[
	P(M_t \ge a) = 2 P(B_t \ge a) = 2 - 2 \Phi(a/\sqrt{t})
\]
thus above it is like we are letting 
\[
	a:= \sup_{0 \le s \le t}B_s-B_t
\]
and 
\[
t:= \sqrt{1-t}
\]
but this seems to me that Naulart is applying a generalization of the reflection principle that is not stated in the text and I could not find a reference for. 
\end{NickQuestion}

\begin{NickQuestion}[Example 6.1.4]
	\begin{enumerate} 
		\item Why does Tanaka's formula imply that $\sup_{x\in \R}\Norm{L_t^x}_p < \infty$?
		
		\item Why does 
		\[
			E(p_\epsilon'(B_s-B_r+B_r-x | \mathcal F_r)) = p_\epsilon'(B_r-x)
		\]
	\end{enumerate}

\end{NickQuestion}

\begin{NickQuestion}[page 91]
	How can we show 
	\[
		\alpha_t = \int_\R \left(\int_0^t \delta_x(B_s)\ud s\right)^2 \ud x = 2 \int_0^t\int_0^v \delta_0(B_v-B_u) \ud u \ud v
	\]
\end{NickQuestion}
\textbf{Solution:}
I believe we have the following 
\begin{align*}
	\int_\R (L_t^x)^2 \ud x & = \int_\R L_t^x L_t^x \ud x
	\\&= \int_0^t L_t^{B_s} \ud s
	\\&= \int_0^t \int_0^t \delta_0(B_v-B_s) \ud v \ud s
	\\& = 2  \int_0^t \int_0^v \delta_0(B_v-B_s) \ud v \ud s
\end{align*}
where I think the last equality is by some sort of symmetry argument and the occupation formula, equation (2.17) in the book, is used in second equality.

\begin{NickQuestion}[Use of backward integral on page 94]
	I believe that there are a series of typos that luckily do not affect the proof. First, $\{B_r-B_s: 0\le s \le r\}$ is not a Brownian motion and I believe he wants to say that $\{\hat{B}_s:=B_r-B_{r-s}: 0\le s \le r\}$ which is a brownian motion. 
	
	Next, Naulart mentions that he is calculating $L_r^{B_r-y}$ and $L_r^{B_r+y}$ for $\hat{B}$ the reason he can do this is because 
	\begin{align*}
		L_r^{B_r-y} &= \int_0^r \delta_0(B_s-B_r+y) \ud s 
		\\&= \int_0^r \delta_0(B_r-B_s -y) \ud s
		\\& = \int_0^r \delta_0(B_r-B_{r-s}-y) \ud s
		\\& = \int_0^r \delta_0(\hat{B}_s - y ) \ud s
		\\& = L_r^y(\hat{B})
	\end{align*}
where $L_r^y(\hat{B})$ is the local time with respect to $\hat{B}$ which we will denote as $L_r^y$. Similarly we can show that $	L_r^{B_r-y} = L_r^{-y}(\hat{B})$ which we will again denote as $L_r^{-y}$.

Now we apply Tanaka's formula for $\hat{B}$ to get the following
\begin{align*}
	L_t^y &= 2(\hat{B}_r-y)_{-}-2(-y)_{-}+2\int_0^r 1_{\{\hat{B}_s < y\}} \ud \hat{B}_s
	\\& = 2(B_r-y)_{-}-2y+2\int_0^r 1_{\{\hat{B}_s < y\}} \ud \hat{B}_s
\end{align*}
where we use the fact that $y>0$ and $\hat{B}_r =B_r$. Now we apply Tanaka's formula again to see
\begin{align*}
L_t^{-y} &= 2(\hat{B}_r+y)_+ - (y)_+ - 2\int_0^r 1_{\{\hat{B}_s > -y\}} \ud \hat{B}_s
\\& = 2(B_r+y)_+ - 2y - 2\int_0^r 1_{\{\hat{B}_s > -y\}} \ud \hat{B}_s
\end{align*}

Therefore 
\begin{align*}
	L_r^{B_r-y} - L_r^{B_r+ y}& = L_r^y - L_r^{-y}
	\\& = 2(B_r-y)_{-}-2y+2\int_0^r 1_{\{\hat{B}_s < y\}} \ud \hat{B}_s - \left(2(B_r+y)_+ - 2y - 2\int_0^r 1_{\{\hat{B}_s > -y\}} \ud \hat{B}_s \right)
	\\& = 2(B_r-y)_{-}-2(B_r+y)_+ + 2\int_0^r 1_{\{|\hat{B}_s| < y\}} \ud \hat{B}_s
	\\&= 2(B_r-y)_{-}-2(B_r+y)_+ + 2\int_0^r 1_{\{|\hat{B}_{r-s}| < y\}} \hat{\ud} B_s
	\\& =  2(B_r-y)_{-}-2(B_r+y)_+ + 2\int_0^r 1_{\{|B_r-B_s| < y\}} \hat{\ud} B_s
\end{align*}
where we use the property of the backward integral that 
\[
	\int_0^T u_s \ud \hat{B}_s = \int_0^T u_{T-s} \hat{\ud}B_s
\]
and also that $\hat{B}_{r-s} = B_r-B_s$. 

Therefore we have that 
\begin{align*}
	\Psi_h(r) &= \int_0^h L_r^{B_r-y} - L_r^{B_r+ y} \ud y
	\\& = \int_0^h L_r^{y} - L_r^{-y} \ud y
	\\& = \int_0^h (2(B_r-y)_{-}-2(B_r+y)_+) \ud y + 2 \int_0^r(h-|B_r-B_s|)_+ \hat{\ud}B_s
\end{align*}
This is different than on page 94 but the $\int \ud y$ integral wont matter because $h \to 0$ and the $\int \hat{\ud}  B_s$ integral is the same. 
\end{NickQuestion}
\begin{NickQuestion}
	Not sure how the conclusion of the proof of (6.8) proves what we want to show. 
\end{NickQuestion}

\begin{NickQuestion}[pg 100] $\;$
	
	Why does 
	\[
		E(\delta_0(B_s^H-B_t^H)) = (2\pi)^{-d/2}|t-s|^{-Hd}
	\]
\end{NickQuestion}
\textbf{Solution for the case when $H=\frac{1}{2}$ and $d=1$ (Brownian Motion)}
\begin{align*}
	E\delta_0(B_s-B_t) &= \lim_{\epsilon \to 0} E\left[\frac{1}{2 \pi \epsilon} \exp\left(-\frac{(B_s-B_t)^2}{2\pi \epsilon}\right)\right]
	\\&= \lim_{\epsilon \to 0}\frac{1}{\sqrt{2\pi(\epsilon + |s-t|)}}
	\\&= \frac{|s-t|^{-1/2}}{\sqrt{2 \pi}}
\end{align*}
where above we just calculate the expectation since $B_s-B_t \sim N(0,s-t)$. I am assuming there is some type of uniform convergence so we can exchange the limit and expected value? 


\section{Chapter 7}
\begin{NickQuestion}[Typos in Prop 7.1.2]
	I am pretty sure there are several typos in the proof of Prop 7.1.2. We are considering all the measurable $g:(-1,1) \to [0,1]$ such that $\int_{-1}^{1}g(y)\ud y = 0$ but this means that $g=0$ almost every where on $(-1,1)$ which does not seem like that is what we want. Also we then consider 
	\[
	\psi(x) = \int_{-\infty}^x g_n(y)\ud y
	\]
	which does not exist and I think needs to be 
	\[
	\psi(x) = \int_{-1}^x g_n(y)\ud y
	\]
	however then the claim is made that 
	\[
		\psi(F) = \int_{-1}^F g_n(y)\ud y \to 0
	\]
	as $n \to \infty$ because $\int_{-1}^1 g(y) \ud y =0$ but this can only happen when $ F \equiv 1$ which isn't always the case. 
\end{NickQuestion}

\begin{NickQuestion}[pg 107]
	Why can we replace $1_{F>x}$ with $1_{F<x}$ in (7.1) on page 106? It says its because the divergence has $0$ expectation but I can't see why that implies the above
	
	We have the following 
	\begin{align*}
	0&= E\left(\delta\left(\frac{DF}{\Norm{DF}_H^2}\right)\right)
	\\&=  E\left((1_{\{F>x\}}+1_{\{F\le x\}})\delta\left(\frac{DF}{\Norm{DF}_H^2}\right)\right)
	\end{align*}
	and therefore we must have 
	\[E\left(1_{\{F>x\}}\delta\left(\frac{DF}{\Norm{DF}_H^2}\right)\right) =-E\left(1_{\{F\le x\}}\delta\left(\frac{DF}{\Norm{DF}_H^2}\right)\right)\]
	which means we cant just replace $1_{\{F>x\}}$ with $1_{\{F<x\}}$ 
\end{NickQuestion}

\begin{NickQuestion}
	\begin{enumerate}
		\item Throughout all the existence of density results, it is claimed that the densities are bounded. Isn't this aways the case? A density needs to satisfy 
	\[
		\int p(x) \ud x = 1
	\]
	and so $p$ must be bounded almost everywhere and if $p$ is continuous then $p$ is bounded everywhere.
	
	\item Also why do we have that the density in Prop 7.2.5 or Theorem 7.3.1 is continuous?
	
		\end{enumerate} 
\end{NickQuestion}
\textbf{solution to 2:} For the case of prop 7.2.5 I believe that 
\[
p(x) = E \left(1_{[F>x]}H_\alpha(F,1)\right)
\]
is continuous because $H_\alpha(F,1) \in \mathbb D_\infty$ and $F$ has components that are all integrable and so I believe that it would be a easy result to show that 
\[
x \mapsto \int_\Omega 1_{[F>x]}(\omega)H_\alpha(F,1)(\omega) \ud P(\omega)
\] 
is a continuous function. This is sort of the same thing as the standard result in real analysis that when we have a function $f \in L^1(\R)$ then 
\[
x \mapsto \int_\R 1_{(-\infty,x]}(t)f(t)\ud t
\]
is continuous.

\begin{NickQuestion}[page 111]
	What is the definition of 
	\[
		\nabla \varphi * \nabla Q_m (x) 
	\]
	since both 	$\nabla \varphi $ and $ \nabla Q_m $ are vectors I do not know that it means to take the convolution between them. 
	
	Also is there a subscript missing on the $\partial \varphi$ in 
	\[
		\sum_{i=1}^m \int_{\R^m} \partial \varphi(x-y) \partial_i Q_m(y) \ud y
	\]
	and if so what should it be? On the next page it says
	\[
		\sum_{i=1}^m \int_{\R^m} \partial_{i} \varphi(x-y) \partial_i Q_m(y) \ud y?
	\]

\end{NickQuestion}

\begin{NickQuestion}
	There are some detail from the proof of Lemma 7.3.2 that are confusing to me. \begin{enumerate}
		\item First why are we assuming $p$ is bounded? Theorem 7.3.1 already says that $p$ is bounded and since that is the case isn't it trivial that there exits a $c$ such that
		\[
		\Norm{p}_\infty \le c \left( \max_{1 \le i \le m} \Norm{H_{(i)}(F,1)}_p \right)^m
		\]
		or is this $c$ above then dependent on $F$ and we do not want that. But either way, what does it mean to assume $p$ to be bounded or not to be bounded? Why would $p$ not be bounded if theorem 7.3.1 already says it is bounded? 
		
		\item Why do we have 
		\[
			\int_{|z-x| \le \epsilon} |z-x|^{(1-m)q}p(x) \ud x \le C_{m,p}\epsilon^{(p-m)/(p-1)}M
		\]
		since they are pulling out the bound $M$ we need to show that 
		\[ 
			\int_{|z-x| \le \epsilon} |z-x|^{(1-m)q} \ud x \le C_{m,p}\epsilon^{(p-m)/(p-1)}
		\]
		but I dont see why this is going to be true because it seems possible that $(m-1)q \ge m$ which would mean that the integral is infinity. But since it seems that we are arbitrarily choosing $p$ and $q$ we can just choose $q\le m/(m-1)$ which avoids this problem. However I still can't figure out how they get $C_{m,p}\epsilon^{(p-m)/(p-1)}$.
	\end{enumerate}
\end{NickQuestion}
\textbf{solution to 2}
By polar coordinates and $q = p/(p-1)$ we can say that 
\begin{align*}
		\int_{|z-x| \le \epsilon} |z-x|^{(1-m)q} \ud x &= \int_{\mathcal S_{m-1}} \int_0^\epsilon r^{\frac{p(1-m)}{p-1}}r^{m-1} \ud r \ud \sigma
		\\&=  \int_{\mathcal S_{m-1}} \int_0^\epsilon r^{\frac{1-m}{p-1}} \ud r \ud \sigma
		\\& = \epsilon^{\frac{p-m}{p-1}}\int_{\mathcal S_{d-1}} \frac{p-1}{p-m} \ud \sigma 
		\\& =  \epsilon^{\frac{p-m}{p-1}} \sigma(\mathcal S_{m-1})\frac{p-1}{p-m}
		\\& = \epsilon^{\frac{p-m}{p-1}} C_{m,p}
\end{align*}
with $C_{m,p}=\sigma(\mathcal S_{m-1})\frac{p-1}{p-m}$

\begin{NickQuestion}[typo on page 114 (covering  vector field)]
	There is a typo on the top of page 114. It says that the covering vector field for $F$ is $u=\gamma_FDF$ but it needs to say $u=\gamma_F^{-1}DF$. For example, in the $d=2$ case we have the following 
	\begin{itemize}
		\item $DF = (DF^1,DF^2)$
		\item $D(\varphi(F)) = \partial_1\varphi(F)DF^1 + \partial_2\varphi(F)DF^2$
		\item $\gamma_F = \left(\langle DF^i, DF^j \rangle\right)_{1 \le i,j \le 2}$
		\item $ \gamma_F^{-1} = \dfrac{1}{|\gamma_F|}
		\begin{pmatrix}
			\langle DF^2, DF^2\rangle\ & -\langle DF^1, DF^2 \rangle\
			\\ -\langle DF^2, DF^1 \rangle\ & \langle DF^1, DF^1 \rangle\
		\end{pmatrix}$
		\item $u := \gamma_F^{-1}(DF^1,DF^2)^T$
	\end{itemize}
then from here it is an easy calculation to see that 
\[
	\langle D(\varphi(F)), u_k \rangle = \partial_k\varphi(F)
\]
\end{NickQuestion}

\begin{NickQuestion}[Possible typo in the book]
	How is itos formula being used at the bottom of page 120?
\end{NickQuestion}
\textbf{possible solution and possible typo in book:} 
It is shown in the book that for 
\[
X_t = X_0 + \int_0^t \sigma(X_s)\ud B_s + \int_0^t b(X_s) \ud s
\]
that we have 
\[
	D_rX_t = \sigma(X_r) + \int_r^t \sigma'(X_s)D_r(X_s) \ud B_s  + \int_r^t b'(X_s)D_r(X_s) \ud s
\]
and note that $D_r(X_r) = \sigma(X_r)$ for each $r \ge 0$. So we can rewrite the above as 
\begin{align} \label{deq}
		D_rX_t = D_r(X_r) + \int_r^t \sigma'(X_s)D_r(X_s) \ud B_s  + \int_r^t b'(X_s)D_r(X_s) \ud s
\end{align}
Now I claim that for
\[
	Y_{r,t} =  \int_r^t \sigma'(X_s) \ud B_s + \int_r^t b'(X_s) - \frac{1}{2}(\sigma')^2(X_s) \ud s 
\]
that $\exp(Y_{r,t})$ solve \eqref{deq}. By Itos formula with the function $f(x)=\exp(x)$ we have that 
\begin{align*}
	\exp(Y_{r,t}) & = \exp(Y_{r,r})  +  \int_r^t \sigma'(X_s)\exp(Y_{r,s}) \ud B_s + \int_0^t \exp(Y_{r,s}) \left(\frac{1}{2}(\sigma')^2(X_s)+b'(X_s) -\frac{1}{2}(\sigma')^2(X_s)\right) \ud s
	\\& = \exp(Y_{r,r})  +  \int_r^t \sigma'(X_s)\exp(Y_{r,s}) \ud B_s + \int_0^t \exp(Y_{r,s}) b'(X_s)\ud s
\end{align*}
and from with we see that $Y_{r,t}$ satisfies \eqref{deq}. However in the text book they say that 
\[
\sigma(X_t) \left( \int_r^t \sigma'(X_s) \ud B_s + \int_r^t b(X_s) - \frac{1}{2}(\sigma')^2(X_s) \ud s \right)
\]
which I believe is incorrect. 
\end{document}
